\paragraph{Data.} We used two public CELLxGENE cohorts: SEA-AD middle temporal gyrus microglia/PVM (10x 3$'$ v3 nuclei only; 84 donors; 34{,}794 nuclei; 18--1{,}068 per donor) \citep{Gabitto2024SEAAD} and Lee et al.\ 2026 fresh prefrontal-cortex microglia \citep{Lee2026}. For Lee we used a subset fixed before this study to keep the file tractable: 80 donors drawn at random, stratified by disease status, from the donors with at least 1{,}600 microglia, and 1{,}600 random cells from each (seed 20261022). Lee is therefore a high-reliability contrast to SEA-AD and over-represents donors with many cells. We analysed \RealG{} genes: the 400 most expressed in both cohorts and 17 further heat-shock, sex-linked and perivascular-macrophage control genes (\RealPairs{} pairs), with identical gene order across cohorts. LCP used $4\times10^5$ permutations on the full cohorts and $10^5$ elsewhere, for all pairs, on GPUs (about 15\,s per full cohort).

\paragraph{Between-donor co-expression calls are common.} With all donors, LCP calls \SEAADfullLCP{} (SEA-AD) and \LEEfullLCP{} (Lee) pairs at FDR 0.05, a fraction of \SEAADfullLCPfrac{} and \LEEfullLCPfrac{} of all pairs; pseudobulk calls \SEAADfullPB{} and \LEEfullPB{}, and the moment-corrected test \SEAADfullMC{} and \LEEfullMC{}. The per-pair reliabilities give $\Deff/D$ medians of \DeffSEAADfull{} and \DeffLEEfull{} at the cohorts' real cell numbers, but only \DeffSEAADcap{} and \DeffLEEcap{} when each donor is reduced to 40 cells (Theorem~\ref{thm:bound}). Sex-linked pairs are a positive control for a purely between-donor association: \emph{XIST} versus Y-linked genes have corrected between-donor correlations of \sexRhoRange{} with within-donor correlations of at most \sexWithinMax{} in absolute value. Heat-shock genes co-vary at both levels (median between-donor correlation \hsRhoSEAAD{} in SEA-AD and \hsRhoLEE{} in Lee; within-donor \hsWSEAAD{} and \hsWLEE{}).

\paragraph{The leak in real data.} At the cohorts' real depths the estimated leak is a small part of the pseudobulk covariance: over the pairs that pseudobulk calls, the median ratio of the estimated leak $D^{-1}\sum_d\kappa_d\hat C_{d,jk}$ to the count-scale covariance of donor means is \SEAADleakFull{} in SEA-AD and \LEEleakFull{} in Lee. This is consistent with Lemma~\ref{lem:leak}, because hundreds of cells per donor make $\bar\kappa$ small. The pairs pseudobulk calls are dominated by strong between-donor associations, however; for pairs that co-vary only within donors the leak is, by definition, the whole signal. To see the leak without reusing cells, we split each donor's cells at random into halves A and B. Half A selected within-donor-only pairs: the top 1\% by within-donor correlation that showed no between-donor evidence (LCP and pseudobulk $p>0.05$; \capNwithin{} pairs). Half B, reduced to 40 cells per donor (three random draws; donors with fewer cells kept all of them), was used only for testing. There the median leak share among pseudobulk calls rises to \SEAADleakSplit{} (SEA-AD) and \LEEleakSplit{} (Lee). In SEA-AD, pseudobulk calls \capPBSEAAD{} of the within-only pairs, uncorrected count-scale means \capnaiveSEAAD{} and LCP \capLCPSEAAD{}; in Lee all methods call \capLCPLEE{}. These counts are small, so the rates are imprecise.

\paragraph{Stability under donor subsampling.} We drew 20 random subsets of 24 and of 40 donors from each cohort and compared their calls with the full-cohort LCP calls, which are themselves estimates. At 24 donors, the fraction of calls that are full-cohort calls of the same sign was \precLCPDtwentyfour{} for LCP, \precPBDtwentyfour{} for pseudobulk and \precMCDtwentyfour{} for the moment-corrected test, which recovered \recMCDtwentyfour{} of the full-cohort calls against \recLCPDtwentyfour{} for LCP. At 40 donors: precision \precLCPDforty{}, \precPBDforty{} and \precMCDforty{}; recall \recLCPDforty{}, \recPBDforty{} and \recMCDforty{}. Because the reference is LCP's own full-cohort call set, this comparison favours LCP and measures stability, not accuracy; the clear finding is that the moment-corrected test recovers few pairs at these donor numbers.

\paragraph{Replication across cohorts.} Of SEA-AD LCP calls, \replLCP{} are LCP calls of the same sign in Lee (pseudobulk: \replPB{}); the rank correlation of pair statistics across cohorts is \rhoLCP{} for LCP and \rhoPB{} for pseudobulk. Crossing the methods gives the same picture: SEA-AD LCP calls are replicated by Lee pseudobulk calls at \xLCPbyPB{}, and SEA-AD pseudobulk calls by Lee LCP calls at \xPBbyLCP{}. The \nPBonly{} pairs called by pseudobulk but not by LCP in SEA-AD replicate as LCP calls in Lee at \replPBonly{}. They are borderline pairs, though, so this lower rate is not by itself evidence of leak. At these cell numbers the leak is small, and the two methods give similar results; the difference matters when donors contribute tens of cells.

\paragraph{Checking the assumptions (post hoc).} Two features of the real data qualify these results. First, depth is not independent of cell state. The pooled within-donor correlation between log library size and a gene's library-normalized count, which is zero in expectation under Lemma~\ref{lem:leak}'s assumption, has median \cplMedSEAAD{} in SEA-AD (median absolute value \cplAbsSEAAD{}) and \cplMedLEE{} in Lee (\cplAbsLEE{}). Its absolute value exceeds \cplStressW{}, the median for the coupled genes of the stress design, for \cplFracSSEAAD{} and \cplFracSLEE{} of the genes, and exceeds 0.1 for \cplFracTSEAAD{} and \cplFracTLEE{}. The within-donor spread of log library size (\sdLibSEAAD{}, \sdLibLEE{}) matches the simulated null (\sdLibNull{}) rather than the stress design (\sdLibStress{}), so the real data are not a replica of the stress design. They do show that depth--state coupling of that size is common, and LCP's stress-design FDR of \STRESSLCPFDR{} applies to the real-data calls as a warning. Second, the donor means have a strong low-rank structure: the first principal component of the standardized donor means carries \pcOneSEAAD{} (SEA-AD) and \pcOneLEE{} (Lee) of their variance, against at most \pcNullSEAAD{} and \pcNullLEE{} after permuting each gene across donors. It is not explained by mean library size, cell number, sex or (SEA-AD) mitochondrial fraction ($|r|\le\pcCorrMaxSEAAD{}$ and $\le\pcCorrMaxLEE{}$). A shared donor factor, biological or technical, is the likely reason that LCP calls \SEAADfullLCPfrac{}--\LEEfullLCPfrac{} of all pairs; adjusting for it changes the estimand and is left for future work.
